% Supplementary manuscript; the evidence manifest remains a separate research artifact.
\documentclass[10pt,twocolumn,letterpaper]{article}
\usepackage[review]{cvpr}
\usepackage{mathtools}
\usepackage{tikz}
\usepackage{pgfplots}
\pgfplotsset{compat=1.18}
\usetikzlibrary{arrows.meta,positioning,fit,calc,backgrounds}
\usepackage{algorithm}
\usepackage[noend]{algpseudocode}
\usepackage{microtype}
\usepackage{array}
\newcommand{\method}{CTA\xspace}
\DeclareMathOperator*{\argmax}{arg\,max}
\DeclareMathOperator*{\argmin}{arg\,min}
\newcommand{\E}{\mathbb{E}}
\newcommand{\R}{\mathbb{R}}
\newcommand{\sg}{\operatorname{sg}}
\newcommand{\bA}{\mathbf{A}}
\newcommand{\cG}{\mathcal{G}}
\newtheorem{proposition}{Proposition}
\AtBeginDocument{\crefname{proposition}{Prop.}{Props.}\Crefname{proposition}{Proposition}{Propositions}%
\crefname{algorithm}{Alg.}{Algs.}\Crefname{algorithm}{Algorithm}{Algorithms}}
\usepackage{placeins}

\definecolor{cvprblue}{rgb}{0.21,0.49,0.74}
\usepackage[pagebackref,breaklinks,colorlinks,allcolors=cvprblue]{hyperref}
\def\paperID{*****}
\def\confName{CVPR}
\def\confYear{2027}
\title{Predict Once, Query Many: World Models that Forecast\\Conditional Trajectory Abstractions\\Supplementary Material}
\author{Anonymous CVPR submission}
\begin{document}
\maketitle
\setcounter{section}{0}
\renewcommand{\thesection}{\Alph{section}}
\section{Implementation details}
This supplement specifies the implemented CTA world model and additional development evidence. Main and supplementary results use reused development roots and one training seed; none is presented as independent multi-seed validation. Privileged actual-future diagnostics below are separate from deployment comparisons.

\paragraph{Visual inputs.}
Frozen DINOv2 ViT-S/14 produces a $16\times16$ feature grid from each resized $224\times224$ image. We project features using a train-fitted 128-dimensional PCA. Context includes the current grid, an $8\times8$ pooled previous grid, and a token embedding $[p_t,p_{t-1}]/512$. The source future includes an endpoint grid, three pooled intermediate grids, and endpoint proprioception. The decoder reconstructs endpoint and intermediate image features, omitting proprioception. Intermediate samples are $(2,4,6)$ for $L=8$ and $(4,8,11)$ for $L=15$; missing samples after native termination repeat the final observation.

\paragraph{Architecture.}
All Transformer components use width 256 and eight heads. The target encoder uses 16 queries and three decoder layers. The reader uses four encoder layers, the feature decoder two decoder layers, and the world model four encoder plus four decoder layers. Its 16 queries predict FSQ coordinate distributions in parallel. Levels $(8,8,4)$ give 256 grid values per source token. The reader receives the $16\times3$ continuous coordinate mean at deployment, with no action input. Neither target encoder nor world model receives the goal.

\paragraph{Representation capacity.}
For a frozen deterministic source encoder,
\begin{equation}
I(\tau;S\mid C)=H(S\mid C)\le16\log_2 256=128\text{ bits}.
\end{equation}
The bound concerns the finite source alphabet, not measured conditional entropy or an entropy-coded deployment rate. Training uses dropout; determinism refers to frozen evaluation mode. The forecast mean consists of 48 float32 values (192 raw bytes), excluding context, goals, weights, and activations. It omits the full predictive distribution. The implemented recipe does not optimize an explicit conditional-entropy loss or a context-only prior.

\paragraph{Task labels.}
For block pose $x$ and target pose $g$, the dense training label is
\begin{equation}
\ell(x,g)=-\frac{1}{512\cdot8}\sum_{v=1}^{8}\|v(x)-v(g)\|_2.
\end{equation}
It uses simulator vertices during training and diagnostic scoring. Deployment uses images, agent positions, proposed actions, and goal features. Native success instead requires post-action coverage greater than .95 before a 300-step episode limit. The reported score is maximum post-action coverage, excluding reset coverage, divided by .95 and clipped to one.

\section{Training and model selection}
Both horizon-specific v2 recipes train from scratch with AdamW: learning rate $3\times10^{-4}$, weight decay .05, warmup 500 updates, dropout .1, and 16 banks/update. Source and predictor training last 16,000 and 10,000 updates, with evaluation every 2,000 updates. Ranking margin is .001, weighted-ranking scale .01, and feature/saturation coefficients are one. Empty eligible-pair sets contribute zero; the weighted ranking denominator is the number of eligible pairs rather than the sum of pair weights.

The feature objective averages normalized endpoint and intermediate MSE. Each normalization uses the training error of copying the current feature grid. A zero-residual copy predictor has expected normalized error one on that distribution. A constant-code decoder can still exploit context and outperform this reference, so feature reconstruction encourages future information without guaranteeing code utilization.

\begin{table}[t]
\centering\small
\begin{tabular}{lrr}
\toprule
Bank source & $L=8$ & $L=15$ \\
\midrule
Policy-state banks & 22,680 & --- \\
On-policy selector states & 10,708 & --- \\
Standard proposals (r4) & 32,590 & 18,606 \\
Perturbed proposals (r4) & 32,590 & 18,606 \\
\midrule
Checkpoint-selection banks & 4,260 & 1,476 \\
\bottomrule
\end{tabular}
\caption{Saved training-bank counts. Siblings share a simulator root and are not independent episodes. Horizon-specific data mixtures differ.}
\end{table}

The training sampler favors banks with useful label spread. For $L=15$, the saved exposure audit reports 57.225\% perturbed banks and 5.557\% mixed native-success banks; these are observed exposure fractions. Training data include states visited by a geometry-assisted selector to broaden support. All branches from an episode remain in one split. Geometry success bonuses and hindsight relabeling are zero in these recipes.

Source training uses separately parameterized CTA codec/reader/decoder, actual-future reader, and original DIR groups in one optimizer with shared global gradient clipping. Their optimization is consequently coupled despite distinct weights. DIR-large is trained separately. Predictor training freezes the target encoder and reader, while gradients through the reader update the world model.

Checkpoint selection maximizes development geometry retained gap, averaging four goal views $(0,5,10,15)$ on a bounded selection subset. Final scoring and acting readers average all sixteen views. All views depict one target pose, without semantic-goal transfer. The selected updates are given below.

\begin{table}[t]
\centering\small
\begin{tabular}{lrr}
\toprule
Selected component & $L=8$ update & $L=15$ update \\
\midrule
Source codec/reader & 16,000 & 8,000 \\
Actual-future reader & 16,000 & 16,000 \\
Original DIR & 14,000 & 12,000 \\
CTA world model & 10,000 & 8,000 \\
END world model & 10,000 & 8,000 \\
\bottomrule
\end{tabular}
\caption{Selected checkpoints from the archived configurations.}
\end{table}

\section{Readout consistency and score stability}
The world-model categorical loss averages cross-entropy over the three coordinates of all 16 tokens. A sum in nats/token is three times this coordinate mean; whole-segment NLL sums over tokens. These loss units are distinct from empirical entropy and source capacity. The reader evaluates the mean code rather than the expected nonlinear score.

\begin{proposition}
For a fixed candidate bank and query, let $u_k$ and $s_k$ be source-reader and predicted-reader scores. Define
\[
\varepsilon=\max_k|(s_k-\bar s)-(u_k-\bar u)|.
\]
For $k_u\in\argmax_k u_k$ and $k_s\in\argmax_k s_k$, $u_{k_u}-u_{k_s}\le2\varepsilon$. If the unique source winner has margin larger than $2\varepsilon$, the predicted winner equals it.
\end{proposition}
\noindent\textit{Proof.} With $\tilde u=u-\bar u$ and $\tilde s=s-\bar s$, centering leaves maximizers unchanged. Then
\begin{align*}
u_{k_u}-u_{k_s}
&=(\tilde u_{k_u}-\tilde s_{k_u})
 +(\tilde s_{k_u}-\tilde s_{k_s})
 +(\tilde s_{k_s}-\tilde u_{k_s})\\
&\le\varepsilon+0+\varepsilon.
\end{align*}
A different predicted winner contradicts this bound when every source margin exceeds $2\varepsilon$.\hfill$\square$

This is a score perturbation bound, rather than a task-performance guarantee. Native-utility regret additionally requires score calibration; closed-loop guarantees require assumptions about dynamics and visited-state distributions.

\section{Additional learned-method results}
\paragraph{Uncertainty.}
The principal $L=8$ ranking estimates use 4,000 root-cluster bootstrap resamples. Retained-gap numerators and denominators are aggregated after root resampling, rather than averaging per-bank ratios. Paired comparisons use identical root resamples. The intervals describe uncertainty over sampled episode roots and do not include variation across trained seeds. All intervals are transcribed from saved summaries.

\begin{table}[t]
\centering\small
\setlength{\tabcolsep}{3pt}
\begin{tabular}{lcc}
\toprule
Difference (pp) & $L=8$ & $L=15$ \\
\midrule
CTA $-$ P0 & +9 [1,17.5] & +9 [1,17] \\
CTA $-$ END & +1 [$-$7.5,9.5] & +6.5 [$-$2.5,15.5] \\
CTA $-$ DIR & +2.5 [$-$5.5,10.5] & +9 [.5,17.5] \\
CTA $-$ DINO-WM & +2 [$-$6,10] & +4.5 [$-$4.5,13.5] \\
\bottomrule
\end{tabular}
\caption{Historical closed-loop paired differences and 95\% root-bootstrap intervals. The same 200 roots are evaluated at each horizon, with different horizon-specific training mixtures.}
\end{table}

\begin{table}[t]
\centering\small
\begin{tabular}{lrr}
\toprule
Exact McNemar $p$ & $L=8$ & $L=15$ \\
\midrule
CTA vs P0 & .044371 & .044371 \\
CTA vs END & .907561 & .182077 \\
CTA vs DIR & .625407 & .050452 \\
CTA vs DINO-WM & .716301 & .385669 \\
\bottomrule
\end{tabular}
\caption{Exploratory binary-success tests without multiplicity correction. The bootstrap interval and exact test at the CTA--DIR $L=15$ boundary do not yield the same nominal threshold decision.}
\end{table}

\paragraph{Candidate-bank expansion.}
The separate $K$-scaling population contains 2,810 P0 decisions from 100 reused development roots. Its newly logged banks differ from the earlier principal $K=8$ experiment: only .4318 of 2,698 common complete $K=8$ banks are identical across those runs, although their initial decision banks agree. Within the scaling log, methods evaluate common banks at each $K$. The table reports mean geometry gains in $10^{-3}$ units; no marginal confidence intervals for these point estimates are implied.

\begin{table}[t]
\centering\small
\begin{tabular}{lrrrr}
\toprule
Mean gain $\times10^3$ & $K=8$ & 16 & 32 & 64 \\
\midrule
CTA WM & .608 & .750 & .877 & .999 \\
END WM & .522 & .644 & .772 & .889 \\
DIR & .424 & .553 & .672 & .745 \\
DIR-large & .437 & .533 & .644 & .739 \\
DINO-WM & .259 & .380 & .454 & .500 \\
\bottomrule
\end{tabular}
\caption{Deployable learned methods on the candidate-expansion log. Point estimates use all decisions in the stated population. Paired root-bootstrap contrasts, rather than marginal error bars, are available.}
\end{table}

At $K=64$, CTA--END gain is $.110\times10^{-3}$ $[.064,.155]\times10^{-3}$; CTA--DIR is $.254\times10^{-3}$ $[.191,.327]\times10^{-3}$; CTA--DIR-large is $.260\times10^{-3}$ $[.172,.347]\times10^{-3}$. An eight-sampled-code readout has lower gain than the implemented mean-code readout by $.306\times10^{-3}$ $[.259,.355]\times10^{-3}$. This is a readout comparison in the tested configuration, not evidence of calibrated predictive uncertainty.

\section{Actual-future diagnostics}
These diagnostics use information unavailable at deployment. FULL scores actual endpoint futures with a separate reader. CODE scores realized source codes from actual future segments. They are not learned acting baselines, and their results are not used as main-paper performance claims. FULL and CODE have different input/readout structures, so their difference is not an isolated compression ablation.

\begin{table}[t]
\centering\small
\setlength{\tabcolsep}{3pt}
\begin{tabular}{lcc}
\toprule
Diagnostic RG & L8 P0 banks & L15 selection banks \\
\midrule
FULL: actual endpoint & .756 [.731,.778] & .776 [.735,.811] \\
CODE: actual code & .628 [.596,.657] & .667 [.608,.719] \\
CTA: predicted code & .680 [.629,.726] & .374 [.265,.478] \\
END: endpoint WM & .635 [.588,.678] & .440 [.313,.548] \\
DIR: action scorer & .483 [.425,.535] & .219 [.071,.349] \\
\bottomrule
\end{tabular}
\caption{Actual-future diagnostics and learned-method references. L8 uses 5,558 P0 decisions / 200 roots; L15 uses 1,476 checkpoint-selection banks / 100 other development roots. The populations and mixtures differ. The columns are not a matched horizon ablation.}
\end{table}

On the L8 bank population, predicted CTA's mean-code input exceeds CODE's realized-code point estimate. The tiers are therefore not a monotone degradation ladder. Smoothing is a possible explanation, without a demonstrated mechanism. On the L15 selection population, CODE exceeds predicted CTA by .294 in point estimate, indicating a substantial forecast/readout discrepancy in that configuration. It is not an independent-test or paired-confidence finding.

\paragraph{Native-success opportunities.}
A crossing bank is one where candidate zero fails within the chunk but some sibling reaches native success. Such banks are uncommon on logged P0 states: 58/5,558 at $L=8$ and 35/3,249 at $L=15$. The diagnostic capture rates below measure selection of a native-success sibling, rather than eventual episode success.

\begin{table}[t]
\centering\small
\begin{tabular}{lrr}
\toprule
Crossing capture & $L=8$ & $L=15$ \\
\midrule
FULL (actual future) & .9655 & .9429 \\
CODE (actual future) & .8966 & .9143 \\
CTA (predicted code) & .7931 & .3714 \\
Geometry selector (privileged) & .9828 & 1.0000 \\
\bottomrule
\end{tabular}
\caption{Capture on small native-crossing subsets. Privileged rows diagnose available distinctions; they are not deployed performance comparators.}
\end{table}

The dense geometry proxy captures nearly all these opportunities. Lower predicted capture consequently points to forecast/readout fidelity, without establishing widespread utility mismatch. Terminated segments repeat final observations, so a native hit is not itself proof of intermediate-event retention. Among 45 L8 P0 roots missing at least one crossing, 33 still eventually succeed, illustrating downstream recovery.

\paragraph{Predictability regularization pilot.}
A separate earlier co-design pilot alternated source and predictor updates. A predictor copy was frozen within each source-update block; it was not a permanently fixed teacher. Raising the predictability penalty from zero to .1 reduced forecast loss while degrading decision distinctions:

\begin{table}[t]
\centering\small
\begin{tabular}{lrr}
\toprule
Earlier pilot & Weight 0 & Weight .1 \\
\midrule
Forecast CE (nats/token) & .558 & .086 \\
Source retained gap & .706 & .458 \\
Token perplexity & 131.2 & 1.37 \\
Sibling distinctness & .765 & .206 \\
\bottomrule
\end{tabular}
\caption{A tested predictability-regularization failure in an earlier model/data regime, separate from the main v2 comparison.}
\end{table}

The paired source-RG change is $-.248$ [$-.297,-.200$]. This supports a collapse diagnosis for that pilot, not a universal conclusion about co-design. Forecast loss alone does not measure the usefulness of its learned target.

\section{Complete component timing}
The archived timing protocol uses CUDA synchronization, ten initial states (roots 2200--2209), three warmups, and ten repeats/state for scorer and context measurements. Policy sampling instead uses one warmup and three repeats/state. Table~\ref{tab:fulltiming} reports medians and p90 values; p90 is an empirical tail summary, not a confidence interval.

\begin{table}[t]
\centering\small
\setlength{\tabcolsep}{3pt}
\begin{tabular}{rlrr}
\toprule
$K$ & Component & $G=1$ & $G=16$ \\
\midrule
8 & CTA WM & 4.11/4.30 & 16.87/17.13 \\
8 & END WM & 4.52/4.64 & 23.40/23.52 \\
8 & DIR & 2.54/2.56 & 26.81/26.83 \\
16 & CTA WM & 4.89/5.15 & 30.52/30.77 \\
16 & END WM & 6.07/6.10 & 44.16/44.28 \\
16 & DIR & 3.94/3.95 & 52.52/52.59 \\
64 & CTA WM & 12.96/13.00 & 114.42/114.70 \\
64 & END WM & 20.09/20.13 & 171.04/171.65 \\
64 & DIR & 13.94/13.97 & 206.67/206.72 \\
\midrule
8 & DINO-WM$^\dagger$ & 28.96/29.15 & 29.07/29.23 \\
16 & DINO-WM$^\dagger$ & 52.27/52.66 & 52.84/53.12 \\
64 & DINO-WM$^\dagger$ & 188.99/189.99 & 190.00/190.52 \\
\bottomrule
\end{tabular}
\caption{\textbf{Median / p90 latency in ms}, H100 3g.40gb MIG. CTA/END/DIR include prediction/readout and exclude shared context and policy. $^\dagger$DINO-WM includes its own encoding, rollout, and cost, with a different precision path. These are component measurements rather than matched total-planner latency.}
\label{tab:fulltiming}
\end{table}

Shared context encoding takes approximately 6.5ms. The common diffusion policy takes 765.33, 777.28, and 923.01ms for $K=8,16,64$. Policy generation dominates this workload. No complete planner-latency or peak-memory result is saved. Our paths use planner mixed precision; equivalent precision with the external DINO-WM path is not established. For $K=8,G=1$, the saved profiler reports 65.92B, 108.50B, and 81.80B operations for CTA, END, and DIR components, respectively; these exclude the shared encoder and policy and depend on profiler operation conventions.

Frozen DINOv2 has 22,056,576 parameters and the proposal policy 262,709,044. The learned counts in the main paper exclude both. DINO-WM predictor/action/proprioception components total 20,122,760 parameters. Reader costs still include context and goal processing, so compact forecast coordinates do not directly measure full inference memory or cost.

\section{Numerical reproducibility}
Historical batched proposals change with grouping. Two L15 evaluations of the same 200 roots and weights report P0 successes 100 and 103, with 43 binary-success flips and 76 step-count mismatches. A privileged geometry selector similarly changes from 152 to 149 successes, with 29 flips and 93 step-count mismatches. Compared source files are byte-identical, while grouping and scorer setup differ. Batch shape is one demonstrated source of sensitivity, without complete attribution of all episode discrepancies. On root 2205, grouping from one to four, 25, or 100 changes initial actions by up to .056335, .037109, or .035614 world units.

Canonical policy sampling fixes denoising microbatch shapes and candidate-prefix construction. A bounded P0 check evaluates root 2205 alone and grouped with root 2206. $K=1,8,16$ action prefixes match exactly; discrete decision fields agree. Four of 128 coverage entries differ by at most $3.33\times10^{-16}$, passing absolute tolerance $10^{-12}$ with zero relative tolerance. This does not qualify every learned scorer, root, simulator state, or hardware backend. The canonical initial bank also differs from its historical bank by up to .017334 world units, so archived success counts cannot replace a fresh canonical evaluation.

Earlier simulator-cloning smoke checks measured live/clone coverage differences near .002 under tolerance .01. Full-state equivalence is not established. Broader restoration, all-arm sampling, and independent-seed control validation remain required. These checks materially limit the interpretation of historical control, while fixed-bank ranking comparisons evaluate the same saved candidates across scorers.
\end{document}
